\appendix
\appendixpage
\addappheadtotoc

\section{Bit length sampler and PMF derivation}
\label{sec:appendix}

\begin{algorithm}[ht]
    \caption{Bit length sampler}
    \label{alg:sampler}
    \SetKwProg{func}{function}{$\colon$}{}
    \KwIn{
        $m$ is the cardinality of the set (or map) to approximate,
        $\fprate$ is the false positive rate,
        $\mu$ is the mean value encoding length.
    }
    \KwOut{
        A random bit length $n$ of a BHF encoding.
    }
    \func{\sampler{$m$, $\fprate$, $\mu$}}{
        Draw a random set $\Set{S}$ of $m$ elements from $\cisb$\;
        \If{$\mu > 0$}{
            Assign random values $v_i$ to each $x_i \in \Set{S}$ with mean
            encoding length $\mu$\;
        }
        $(h_k, b) \gets \MakeBHF(\Set{S}, \fprate, 0)$\;
        \Return $\BL(h_k) + \BL(b) + \mathcal{O}(1)$\;
    }
\end{algorithm}

\begin{theorem}[PMF of the random bit length]
\label{thm:pmf_detailed}
The random bit length $\RV{L}$ of the salt $b$ has the probability mass
function
\begin{equation}
    \PDF{n \Given m, \fprate, \mu}[\RV{L}]
    = q^{2^n - 1}\!\left(1 - q^{2^n}\right)\,,
\end{equation}
where $q = 1 - p$ and $p$ is the per-trial success probability.
\end{theorem}
\begin{proof}
Each candidate salt is an independent Bernoulli trial with success probability
$p$.
The algorithm tests candidates in order of increasing bit length, with $2^n$
candidates of length $n$.

For $\RV{L} = n$ to occur:
\begin{enumerate}
    \item All $2^n - 1$ candidates of length $< n$ must fail.
    Each fails independently with probability $q = 1 - p$, so the joint
    failure probability is $q^{2^n - 1}$.

    \item At least one candidate of length $n$ must succeed.
    All $2^n$ candidates of length $n$ fail with probability $q^{2^n}$, so
    the complementary probability is $1 - q^{2^n}$.
\end{enumerate}
By independence, $\Prob{\RV{L} = n} = q^{2^n - 1}(1 - q^{2^n})$.

\paragraph{Verification that this is a valid PMF.}
The sum telescopes:
\begin{align}
    \sum_{n=0}^{\infty} q^{2^n - 1}(1 - q^{2^n})
    &= \sum_{n=0}^{\infty} \left(q^{2^n - 1} - q^{2^{n+1} - 1}\right) \\
    &= (1 - q) + (q - q^3) + (q^3 - q^7) + (q^7 - q^{15}) + \cdots \\
    &= 1\,.
\end{align}

\paragraph{Alternative derivation.}
Since $\RV{L} = \lfloor \log_2 \RV{Q} \rfloor$ where
$\RV{Q} \sim \geodist(p)$:
\begin{align}
    \Prob{\RV{L} = n}
    &= \Prob{2^n \leq \RV{Q} < 2^{n+1}} \\
    &= \sum_{j=2^n}^{2^{n+1}-1} p(1-p)^{j-1} \\
    &= p \cdot \frac{q^{2^n - 1} - q^{2^{n+1} - 1}}{1 - q}
    = q^{2^n - 1}(1 - q^{2^n})\,.
\end{align}
\end{proof}

\begin{theorem}[Expected encoding size]
\label{thm:expected_size}
The expected salt bit length is
\begin{equation}
    \Expect{\RV{L}} = \sum_{n=1}^{\infty} q^{2^n - 1}
    = q + q^3 + q^7 + q^{15} + \cdots\,,
\end{equation}
where $q = 1 - p$.
\end{theorem}
\begin{proof}
\begin{align}
    \Expect{\RV{L}}
    &= \sum_{n=0}^{\infty} n\, q^{2^n - 1}(1 - q^{2^n}) \\
    &= \sum_{n=0}^{\infty} n\left(q^{2^n - 1} - q^{2^{n+1} - 1}\right)\,.
\end{align}
Expanding the first few terms:
\begin{align}
    \Expect{\RV{L}}
    &= 0 \cdot (1 - q) + 1 \cdot (q - q^3)
        + 2 \cdot (q^3 - q^7) + 3 \cdot (q^7 - q^{15}) + \cdots \\
    &= q + q^3 + q^7 + q^{15} + \cdots
    = \sum_{n=1}^{\infty} q^{2^n - 1}\,.
\end{align}
\end{proof}

\section{Adaptive threshold PMF derivation}
\label{sec:adaptive_pmf_derivation}

We derive the law of the threshold $T = \OS{p}{m}$ stated in
\cref{thm:fpr_distribution} and relate it to the distinct-residue formula
printed in earlier versions of this paper.

\begin{theorem}[Law of the $p$-th order statistic of discrete uniform draws]
\label{thm:os_pmf_detailed}
Let $R_1,\ldots,R_m$ be i.i.d.\ $\operatorname{Uniform}\{0,1,\ldots,N-1\}$.
Then for $j = 0, 1, \ldots, N-1$,
\begin{equation}
\label{eq:os_pmf_detailed}
    \Prob{\OS{p}{m} \leq j}
    \;=\; \sum_{a=p}^{m} \binom{m}{a}
    \left(\frac{j+1}{N}\right)^{a}\left(1 - \frac{j+1}{N}\right)^{m-a}\,,
\end{equation}
and $\Prob{\OS{p}{m} = j} = \Prob{\OS{p}{m} \leq j} - \Prob{\OS{p}{m} \leq j-1}$.
\end{theorem}
\begin{proof}
Let $A_j = \Card{\{i : R_i \leq j\}}$.
The $p$-th smallest residue is at most $j$ if and only if $A_j \geq p$.
The indicators $\mathbf{1}[R_i \leq j]$ are independent with success
probability $(j+1)/N$, so $A_j \sim \bindist(m, (j+1)/N)$, which gives
\cref{eq:os_pmf_detailed}.
\end{proof}

\paragraph{Distinct-residue approximation.}
Earlier versions stated, as the exact law for i.i.d.\ residues, the formula
\begin{equation}
\label{eq:os_pmf_distinct}
    \Prob{\OS{p}{m} = j \mid D}
    \;=\; \frac{\binom{j}{p-1}\binom{N-1-j}{m-p}}{\binom{N-1}{m-1}}
    \cdot \frac{m}{N}\,,
    \quad j = p-1, \ldots, N-1-(m-p)\,,
\end{equation}
which is the law for $m$ residues drawn \emph{without replacement}.
It equals the law of $\OS{p}{m}$ conditioned on the event $D$ that all $m$
residues are distinct, since given $D$ the residues are a uniformly random
$m$-subset.
It is not the unconditional law: for $N = m = 2$ and $p = 1$,
\cref{eq:os_pmf_distinct} gives $\Prob{T = 0} = 1$, whereas
\cref{eq:os_pmf_detailed} and direct enumeration give
$(\Prob{T=0}, \Prob{T=1}) = (3/4, 1/4)$.
The two laws differ in total variation by at most
$\Prob{D^c} \leq m(m-1)/(2N)$, because for any event $E$,
$\Prob{E} - \Prob{E \mid D} = \Prob{E \cap D^c} - \Prob{E \mid D}\Prob{D^c}$,
which lies in $[-\Prob{D^c}, \Prob{D^c}]$.
So \cref{eq:os_pmf_distinct} is a birthday-scale approximation, useful only
when $m^2 \ll N$.
An earlier claim that $T+1$ follows a shifted beta-binomial law is also
withdrawn.
The exact law was checked against exhaustive enumeration with exact rational
arithmetic for all $N \leq 5$, $m \leq 4$, $1 \leq p \leq m$, and the bound for
all of these cases with $m \leq N$.

\paragraph{Continuous limit.}
By the standard identity
$\Prob{\bindist(m, x) \geq p} = I_x(p, m - p + 1)$~\cite{order_statistics},
\cref{eq:os_pmf_detailed} reads $\Prob{\OS{p}{m} \leq j} = I_{(j+1)/N}(p, m-p+1)$
exactly.
Setting $x = (j+1)/N$ and letting $N \to \infty$ with $m$ fixed, the law of
$(\OS{p}{m}+1)/N$ converges to the $\betadist(p, m - p + 1)$ distribution,
confirming \cref{eq:fpr_beta}.
